\hypertarget{cont9:paper-ix-compatible-towers-with-prescribed-parameters}{%
\subsection{Chapter IX: compatible towers with prescribed
parameters}\label{cont9:paper-ix-compatible-towers-with-prescribed-parameters}}

The construction gives actual compatible polynomial towers from initial
solved-block surjectivity, for every compatible integer parameter path.
For a prime base, it constructs such a path for every p-adic parameter
in the initial ball and proves an exact polynomial endpoint with all
retained precisions. The formal companion identifies the six supporting modules and their
recorded compiler and axiom-audit results.

\hypertarget{cont9:statement-and-notation}{%
\subsubsection{Statement and notation}\label{cont9:statement-and-notation}}

Let \(q>0\), \(L\ge1\), and \[
F\in\mathbb Z[Y_1,\ldots,Y_m,T_1,\ldots,T_k]^e,
\qquad x=(y_0,t_0)\in\mathbb Z^{m+k},
\qquad F(x)\in q^L\mathbb Z^e.
\] Write \(S_q=\mathbb Z/q\mathbb Z\), and let the actual reduced
Jacobian blocks be \[
J_x=D_YF(x)\bmod q:S_q^m\longrightarrow S_q^e,
\qquad K_x=D_TF(x)\bmod q:S_q^k\longrightarrow S_q^e.
\] Assume \(J_x\) is surjective. A vector divisibility condition below
is coordinatewise.

\textbf{Integer prescribed-path theorem.} For every supplied path
\(t:\mathbb N\to\mathbb Z^k\) satisfying \[
t(0)=t_0,\qquad t(n+1)-t(n)\in q^{L+n}\mathbb Z^k,
\] there is a single tower \(z:\mathbb N\to\mathbb Z^{m+k}\) with \[
z(0)=x,\qquad \operatorname{free}(z(n))=t(n),
\] \[
F(z(n))\in q^{L+n}\mathbb Z^e,\qquad
z(n+1)-z(n)\in q^{L+n}\mathbb Z^{m+k}.
\] The parameter equality is exact in integers. For \(a\le b\), \[
z(b)-z(a)\in q^{L+a}\mathbb Z^{m+k},\qquad
DF(z(n))\bmod q=DF(x)\bmod q.
\] Consequently \(J_{z(n)}=J_x\) at every level. No later surjectivity
or per-step solvability is assumed. This theorem includes positive
composite bases and \(q=1\).

\textbf{Arbitrary p-adic target theorem.} Now let \(q=p\) be prime. For
every \[
\tau\in t_0+p^L\mathbb Z_p^k,
\] there are a compatible integer tower \(z\) as above and
\(b\in\mathbb Z_p^{m+k}\) such that \[
z(n)\longrightarrow b,\qquad F(b)=0,\qquad
\operatorname{free}(b)=\tau,
\qquad b-z(n)\in p^{L+n}\mathbb Z_p^{m+k}\quad(n\ge0).
\] A compatible integer parameter path converging to \(\tau\) is
constructed in the proof, with its initial representative fixed exactly.
A separate theorem accepts every supplied compatible integer path with
limit \(\tau\) and retains its exact coordinates in the constructed
tower.

\textbf{Exact branching.} At any stage \(n\), put \(N=L+n\). For every
parameter digit \(v\in\mathbb F_p^k\), let \[
\mathcal L_{n,v}=\left\{\bar w\in(\mathbb Z/p^{N+1}\mathbb Z)^{m+k}:
\begin{array}{l}
\bar w\equiv z(n)\pmod{p^N},\\
F(\bar w)=0,\\
\operatorname{free}(\bar w)=t(n)+p^Nv\pmod{p^{N+1}}
\end{array}\right\}.
\] Here \(p^Nv\) means any integer representative of \(v\), multiplied
by \(p^N\). The residue class does not depend on the representative.
Then \[
\#\mathcal L_{n,v}=p^{m-e}\qquad\text{for every }n\text{ and every }v.
\] Surjectivity implies \(e\le m\). This is a count of distinct actual
next-level residue classes. It does not count integer representatives or
entire infinite towers.

\textbf{Square solved-block uniqueness.} If \(m=e\), every prescribed
\(\tau\) has exactly one root in the original ball: \[
\exists!\,b\in x+p^L\mathbb Z_p^{e+k}:\quad
F(b)=0,\qquad\operatorname{free}(b)=\tau.
\] Thus all valid integer parameter approximations and all lifting
choices with the same target have the same endpoint. This does not
assert uniqueness of the integer representatives in the tower. All prime
statements include \(p=2\).

\hypertarget{cont9:proof}{%
\subsubsection{Proof}\label{cont9:proof}}

At depth \(N\ge1\), suppose \(F(a)=q^Nr\) and a desired next integer
parameter target satisfies \(t'-\operatorname{free}(a)=q^Nw\), with
\(w\in\mathbb Z^k\). Surjectivity supplies a solved digit \(u\in S_q^m\)
satisfying \[
J_a u=-\bar r-K_a\bar w.
\] Choose an integer representative \(\widetilde u\), and set \[
a'=a+q^N(\widetilde u,w).
\] The parameter quotient is the supplied integer \(w\). Therefore
\(\operatorname{free}(a')=t'\) exactly. The checked polynomial Taylor
identity gives \[
F(a')\equiv F(a)+q^NDF(a)(\widetilde u,w)
\equiv0\pmod{q^{N+1}}.
\] Its quadratic remainder is divisible by \(q^{2N}\), and
\(2N\ge N+1\). Also \(a'\equiv a\pmod q\), so evaluation of every
polynomial partial derivative gives the same reduced Jacobian. Induction
on the level, carrying the root, exact parameters, and fixed solved
Jacobian as its state, constructs the entire integer tower. The
existence proof uses ordinary classical choice. It does not supply an
optimized computational lifting algorithm.

For a p-adic target \(\tau\), use \(t(0)=t_0\) and, for \(n\ge1\), the
canonical integer representatives of \(\tau\bmod p^{L+n}\). The initial
congruence hypothesis and the truncation precision imply \[
\tau-t(n)\in p^{L+n}\mathbb Z_p^k.
\] Subtracting successive precision relations gives compatibility in
\(\mathbb Z_p\). The checked faithful integer-divisibility theorem gives
compatibility in \(\mathbb Z\). These precision relations also give
convergence to \(\tau\).

The compatible vector-tower completion theorem constructs \(b\), proves
convergence, and retains every stage precision. Continuous polynomial
evaluation and the increasing root precision give \(F(b)=0\). Each
parameter coordinate has both the corresponding coordinate of \(b\) and
\(\tau\) as its limit, so uniqueness of limits gives exact parameter
equality.

For every proposed next parameter digit, the previously checked
obstruction is \([\bar r+K_av]\in S_p^e/\operatorname{im}J_a\). The
fixed solved Jacobian is surjective, so this class vanishes for every
digit at every stage. The actual lift fiber is a torsor for
\(\ker J_a\). Rank-nullity gives its cardinality \(p^{m-e}\).

When \(m=e\), the new solved-Jacobian definition is definitionally the
existing square block used by the checked continuous parameter chart.
Every root in the initial ball with the given parameters is that chart
point. Applying this to any two roots proves equality, including
endpoints obtained by different constructions.

\hypertarget{cont9:scope-controls}{%
\subsubsection{Scope controls}\label{cont9:scope-controls}}

The inherited checked example \(F(Y,T)=T\), \(p=2\), \(x=0\), \(L=1\)
has a surjective full Jacobian but zero solved block. Prescribing
parameter digit \(1\) forces \(T=2\pmod4\), so no next root exists.
Full-Jacobian surjectivity cannot replace the hypothesis on the
prescribed solved block.

The new checked example \(F(Y_1,Y_2,T)=Y_1\) has a surjective solved
block \((1,0)\). Nevertheless \((0,0,0)\) and \((0,2,0)\) are distinct
exact roots in the same modulo-two ball with identical parameters.
Uniqueness therefore requires the square solved case. The corresponding
negative control deliberately attempts the stronger false conclusion and
must be rejected by Lean.

Depth \(L\ge1\) is essential for the Taylor step used here. The general
singular case can have a nonzero next obstruction after an earlier zero
obstruction. Neither the finite one-step criterion nor this smooth
theorem establishes unrestricted singular persistence.
